\section{Purpose, scope, and main result}
\label{sec:purpose}

The purpose of this note is to apply Algorithm~1 completely to a networked
SEIR model with three subpopulations. The example is deliberately small enough
for every retained observable and every newly generated nonlinear term to be
listed explicitly, while still exhibiting all features of the general
construction:

\begin{itemize}[leftmargin=2em]
\item diagonal channels representing transmission within each city;
\item off-diagonal channels representing transmission between cities;
\item a mechanistic quadratic seed set;
\item linear closure within each homogeneous polynomial degree;
\item generation of the next degree by the quadratic transmission field;
\item commutative canonicalization and duplicate merging;
\item an exact block hierarchy and a next-degree truncation residual.
\end{itemize}

The complete calculation gives
\[
|\seedM_2^{\mathcalE}|=\NumChannels,
\qquad
|\mathcal M_2^{\mathcalE}|=\NumMtwo,
\]
\[
|\seedM_3^{\mathcalE}|=\NumMthreeSeeds,
\qquad
|\mathcal M_3^{\mathcalE}|=\NumMthree,
\]
and therefore
\[
\dim z_2^{\mathcalE}=\DimZtwo,
\qquad
\dim z_3^{\mathcalE}=\DimZthree.
\]
The \NumMthreeClosure{} cubic coordinates added by linear closure are essential:
they ensure that the cubic dictionary is invariant under the degree-preserving
linear SEIR transitions. Once this closure has been enforced, the residual of
the order-three truncation contains only quartic terms.

The document has three complementary roles. It is:

\begin{enumerate}[leftmargin=2em]
\item a mathematical derivation of the example;
\item a traceability record for the dimension counts used in the article;
\item a reproducibility check against the accompanying Python script and CSV
      files.
\end{enumerate}
